\chapter*{Abstract}
\addcontentsline{toc}{chapter}{Abstract}



The rapid proliferation of Twenty20 (T20) cricket as the dominant format of international competition has created a growing demand for accurate, real-time predictive tools that can assess scoring trajectories from incomplete innings data. Despite considerable prior work employing conventional statistical and machine learning methods, the potential of Large Language Models (LLMs) for domain-specific sports forecasting tasks remains largely unexplored. This thesis investigates whether a compact, open-source LLM, enhanced through parameter-efficient fine-tuning, can serve as an effective in-game score predictor for T20 cricket first innings.

A dataset comprising international men's T20 matches sourced from Cricsheet, spanning the period 2005 to 2026, was processed into discrete in-game state snapshots taken at two-over intervals. Each snapshot was characterised by fourteen contextual features, including cumulative runs scored, wickets lost, current run rate, dot ball counts, boundary frequencies, powerplay completion status, and a venue par score derived from historical match outcomes. A novel feature, the \textit{Remaining Batting Strength Score}, was engineered as a weighted composite of the player types yet to bat, providing a quantitative encoding of residual scoring capacity. These structured snapshots were serialised into natural-language prompts and paired with the actual final first-innings total to form a labelled training corpus. A strictly temporal train--validation--test partitioning strategy was adopted to prevent data leakage and ensure evaluation fidelity under realistic deployment conditions.

Three experimental conditions were evaluated on a held-out test set using Mean Absolute Error (MAE), Mean Squared Error (MSE), and the coefficient of determination ($R^2$) as performance metrics. First, the base \textit{Llama 3.2-3B} model was assessed in a zero-shot setting to establish a pre-training baseline. Second, the same model was subjected to Supervised Fine-Tuning (SFT) using Quantised Low-Rank Adaptation (QLoRA), training exclusively on the domain-specific prompt--completion dataset. Third, the frontier model \textit{GPT-4.1-nano} was evaluated zero-shot to provide an upper-bound reference from a substantially larger, instruction-tuned proprietary system.

The findings of this study demonstrate that domain-specific fine-tuning of a small, resource-efficient LLM yields substantial improvements over the zero-shot baseline and provides a competitive alternative to proprietary frontier models for the task of in-game T20 score prediction. The work contributes a reusable data pipeline, a reproducible fine-tuning framework, and empirical evidence supporting the viability of small fine-tuned LLMs as practical sports analytics tools, particularly in resource-constrained settings.
